\documentclass[11pt]{article}
\usepackage[letterpaper,margin=0.85in]{geometry}\usepackage[T1]{fontenc}\usepackage[utf8]{inputenc}\usepackage{amsmath,amssymb,enumitem,booktabs,array,xcolor,fancyhdr,tabularx}
\setlist{itemsep=2pt,topsep=3pt}\pagestyle{fancy}\fancyhf{}\renewcommand{\headrulewidth}{0.3pt}
\fancyhead[L]{\small 4.4 Part 2 $\cdot$ Related Rates: Differentiate in Time}\fancyhead[R]{\small Teacher key}\fancyfoot[C]{\small\thepage}
\definecolor{ans}{RGB}{28,125,74}\newcommand{\A}[1]{\textcolor{ans}{\textbf{#1}}}
\begin{document}

\section*{4.4 Part 2 $\cdot$ Related Rates: Differentiate in Time --- teacher key}
\textbf{Question:} How do I turn a link equation into an equation about rates?\par
\subsection*{Self-paced lesson}No clock, no trimming, no skipping, no end mark, no stamp. Every tab is required in order; progress saves on the device so a student can stop and finish the next day.\par
\subsection*{Warm-up review (3 easy items)}\begin{itemize}
\item The derivative of $x^2$ with respect to $x$ is... \A{$2x$.} Power rule: bring the 2 down and lower the power by one.
\item $\dfrac{d}{dt}\bigl(r^3\bigr)$, where $r$ depends on $t$, equals... \A{$3r^2\dfrac{dr}{dt}$.} Chain rule: differentiate the power, then multiply by the rate of $r$.
\item The link is $A=\pi r^2$. The number that should be put in last is... \A{$r=10$ (true for one moment).} Numbers true for one moment wait until the last move.
\end{itemize}
\subsection*{Independent lesson}No huddle, no teacher talk, no codes. Students read a short picture tutorial (Concept), watch worked problems move by move (I do), then solve generated problems (You do). Add \texttt{\#key} to the lesson address to open every tab with answers marked.\par
\subsection*{Targets}\begin{itemize}\item I can differentiate a link equation with respect to time.
\item I can put the moment's numbers in last and solve for the wanted rate.
\item I can say what the answer means, with its sign and units.\end{itemize}
\subsection*{Predictions}\begin{itemize}
\item A puddle's radius grows 2 cm every second. When the radius is 10 cm, the area grows faster than when the radius is 5 cm. \A{True: a bigger circle gains more area each second.} The area's rate depends on $r$ itself. A longer edge adds more area each second.
\item The radius grows 3 cm/s. The area grows $\dfrac{dA}{dr}=20$ cm\textsuperscript{2} per cm of radius. The area grows... \A{60 cm\textsuperscript{2} per second.} Rates multiply along the chain: $20\times3=60$.
\end{itemize}
\section*{Tutorial script, with every question and answer}
\subsection*{Concept tutorial (7 pictures)}
\textit{I can say why a rate rides along when I differentiate in time.}\par
\begin{enumerate}
\item \textbf{Time leads to the radius.} As time $t$ passes, the radius $r$ changes. That speed is $\dfrac{dr}{dt}$. \textit{(Picture: chain. Three boxes: time $t$, radius $r$ and area $A$. The first red arrow says how fast $r$ changes as time passes.)}\\ \textit{Question:} Which symbol is the speed of the radius? \A{Answer: $\dfrac{dr}{dt}$.} The speed of $r$ in time is $\dfrac{dr}{dt}$.
\item \textbf{The radius leads to the area.} A change in $r$ makes $A$ change too. How much $A$ changes per unit of $r$ is $\dfrac{dA}{dr}$. \textit{(Picture: chain. The second red arrow says how fast $A$ changes per unit of $r$.)}\\ \textit{Question:} What does $\dfrac{dA}{dr}=12$ cm\textsuperscript{2}/cm say? \A{Answer: Each extra cm of radius adds about 12 cm\textsuperscript{2} of area..} It compares area to radius, with no time in it.
\item \textbf{Multiply along the chain.} Going from $t$ all the way to $A$, we multiply the rates on the way. \textit{(Picture: chain. The green path goes from $t$ to $A$ the long way. The rates on the way are multiplied.)}\\ \textit{Question:} $\dfrac{dA}{dr}=12$ cm\textsuperscript{2}/cm and $\dfrac{dr}{dt}=0.5$ cm/s. Find $\dfrac{dA}{dt}$. \A{Answer: 6 cm\textsuperscript{2}/s.} $12\times0.5=6$ cm\textsuperscript{2}/s.
\item \textbf{A squared quantity is steeper.} The line $r$ has the same slope all the time. The curve $r^2$ gets steeper as $r$ gets bigger. \textit{(Picture: chaingraph. Left graph: $r$ against time, a straight line of slope $0.5$. Right graph: $r^2$ against time. The red lines are slopes. Slide the time.)}\\ \textit{Question:} At $r=2$ with $\dfrac{dr}{dt}=0.5$, the slope of $r^2$ is $2r\cdot\dfrac{dr}{dt}$. What is it? \A{Answer: 2.} $2\times2\times0.5=2$. Look at the readout when $t=2$.
\item \textbf{The chain factor.} For $r^3$ the slope is $3r^2\cdot\dfrac{dr}{dt}$. The factor $\dfrac{dr}{dt}$ always rides along. \textit{(Picture: chaingraph. The right graph is now $r^3$. Its slope at each time is shown in the readout line.)}\\ \textit{Question:} Which is $\dfrac{d}{dt}(r^3)$? \A{Answer: $3r^2\dfrac{dr}{dt}$.} Power rule gives $3r^2$, then multiply by the rate of $r$.
\item \textbf{Fixed numbers have no rate.} The radius $R$ is a fixed number, so it has no rate. Only $h$ brings a rate. \textit{(Picture: cylinder. A tank with a fixed radius $R$. Only the depth $h$ changes. The slider moves $h$.)}\\ \textit{Question:} For $V=\pi R^2h$ with $R$ fixed, $\dfrac{dV}{dt}$ is... \A{Answer: $\pi R^2\dfrac{dh}{dt}$.} $\pi R^2$ is a fixed number multiplying $h$, so only $\dfrac{dh}{dt}$ appears.
\item \textbf{Now put the moment in.} After differentiating, we put the moment's numbers in. Here $\dfrac{dA}{dt}=2\pi r\dfrac{dr}{dt}$ with $r=10$ and $\dfrac{dr}{dt}=2$. \textit{(Picture: balloon. The puddle at the moment $r=10$ cm.)}\\ \textit{Question:} Find $\dfrac{dA}{dt}=2\pi r\dfrac{dr}{dt}$ when $r=10$ and $\dfrac{dr}{dt}=2$. Give it to 2 decimal places. \A{Answer: 125.66 cm\textsuperscript{2}/s.} $2\pi(10)(2)=40\pi\approx125.66$ cm\textsuperscript{2}/s.
\end{enumerate}
\subsection*{I do: the five moves}\begin{enumerate}\item Draw and name
\item List what you know and want
\item Link with one equation
\item Differentiate in t
\item Plug in last\end{enumerate}
\subsection*{I do (8 steps)}
\textit{I can watch moves 4 and 5 done on two problems and spot the common slip.}\par
\begin{enumerate}
\item \textbf{Moves 1 to 3 are done.} \textit{Moves 1 to 3 $\cdot$ Name, list, link.} We did these in Part 1. Here is what we have. \textit{(Picture: balloon. The puddle from above. The red arrows show the radius growing.)}\\ Given: $\dfrac{dr}{dt}=2$ cm/s\\ Want: $\dfrac{dA}{dt}$ when $r=10$ cm\\ Link: $A=\pi r^2$\\ \textit{Question:} Which number is saved for the LAST move? \A{Answer: $r=10$.} $r=10$ is true for one moment only, so it waits.
\item \textbf{Differentiate both sides.} \textit{Move 4 $\cdot$ Differentiate in t.} Take $d/dt$ of both sides: $A$ gives $\dfrac{dA}{dt}$, and $\pi r^2$ gives $2\pi r\dfrac{dr}{dt}$. \textit{(Picture: chain. The chain from $t$ to $r$ to $A$. Rates are multiplied along it.)}\\ $\dfrac{dA}{dt}=2\pi r\,\dfrac{dr}{dt}$\\ \textit{Question:} Which line is the derivative in time? \A{Answer: $\dfrac{dA}{dt}=2\pi r\dfrac{dr}{dt}$.} Power rule gives $2\pi r$, and the chain factor $\dfrac{dr}{dt}$ rides along.
\item \textbf{Plug in the moment.} \textit{Move 5 $\cdot$ Plug in last.} Now put in $r=10$ and $\dfrac{dr}{dt}=2$. Then multiply. \textit{(Picture: balloon. The puddle at the moment $r=10$ cm.)}\\ $\dfrac{dA}{dt}=2\pi(10)(2)$\\ $\dfrac{dA}{dt}=40\pi\approx125.7$ cm\textsuperscript{2}/s\\ \textit{Question:} What does 125.7 cm\textsuperscript{2}/s tell us? \A{Answer: The area is growing 125.7 cm\textsuperscript{2} each second at that moment..} Positive means growing, and cm\textsuperscript{2}/s is area per second.
\item \textbf{Problem 2: the sliding ladder.} \textit{Moves 1 to 3 $\cdot$ Name, list, link.} \textbf{Problem:} A 13 ft ladder slides away from a wall. The foot moves out at 2 ft/s. How fast is the top falling when the foot is 5 ft from the wall? We name, list and link as before. The 13 stays because the ladder never changes length. \textit{(Picture: ladder. A 13 ft ladder against a wall. $x$ is the foot, $y$ is the top.)}\\ Given: $\dfrac{dx}{dt}=2$ ft/s\\ Want: $\dfrac{dy}{dt}$ when $x=5$\\ Link: $x^2+y^2=13^2$\\ \textit{Question:} Which number of the 13 ft ladder problem is a FIXED number? \A{Answer: 13, the ladder length.} Only the ladder length never changes.
\item \textbf{Differentiate both sides.} \textit{Move 4 $\cdot$ Differentiate in t.} The 13 is fixed, so $13^2$ has rate 0. Each of $x^2$ and $y^2$ brings its own rate. \textit{(Picture: ladder. The red arrows show $x$ growing and $y$ shrinking.)}\\ $2x\dfrac{dx}{dt}+2y\dfrac{dy}{dt}=0$\\ \textit{Question:} Why is the right side $0$? \A{Answer: $13^2$ is a fixed number, so its rate is 0..} A fixed number has no rate.
\item \textbf{Find the missing height first.} \textit{Move 5 $\cdot$ Plug in last.} We need $y$ at this moment. Use the link equation with $x=5$: $25+y^2=169$. \textit{(Picture: ladder. The ladder at the moment $x=5$. The top height $y$ is not given.)}\\ At $x=5$: $y^2=169-25=144$, so $y=12$\\ \textit{Question:} At $x=5$ on the 13 ft ladder, how high is the top? \A{Answer: 12 ft.} $y^2=169-25=144$, so $y=12$ ft.
\item \textbf{Plug in and solve.} \textit{Move 5 $\cdot$ Plug in last.} Put in $x=5$, $y=12$ and $\dfrac{dx}{dt}=2$. Solve for $\dfrac{dy}{dt}$. \textit{(Picture: ladder. All the moment's numbers are now filled in.)}\\ $2(5)(2)+2(12)\dfrac{dy}{dt}=0$\\ $20+24\dfrac{dy}{dt}=0$\\ $\dfrac{dy}{dt}=-\dfrac{20}{24}=-\dfrac56\approx-0.83$ ft/s\\ \textit{Question:} Solve $20+24\dfrac{dy}{dt}=0$ for $\dfrac{dy}{dt}$. Give it to 2 decimal places. \A{Answer: -0.83 ft/s.} $24\dfrac{dy}{dt}=-20$, so $\dfrac{dy}{dt}=-0.83$ ft/s. The top is falling, so the rate is negative.
\item \textbf{The common slip.} Slip: writing $\dfrac{dA}{dt}=2\pi r$ and stopping. The chain factor $\dfrac{dr}{dt}$ was dropped, so the answer would have the wrong units. \textit{(Picture: chain. The chain is missing its third step, which is the green path that multiplies by $\dfrac{dr}{dt}$.)}\\ Wrong: $\dfrac{dA}{dt}=2\pi r$\\ Right: $\dfrac{dA}{dt}=2\pi r\dfrac{dr}{dt}$\\ \textit{Common slip.}\\ \textit{Question:} How can you tell the slip from the units? \A{Answer: $2\pi r$ is in cm, but $\dfrac{dA}{dt}$ must be cm\textsuperscript{2} per second..} cm is a length. A rate of area needs cm\textsuperscript{2}/s. The missing $\dfrac{dr}{dt}$ brings the "per second".
\end{enumerate}
\subsection*{Big ideas}\begin{itemize}\item Differentiate both sides of the link with respect to $t$. Every changing letter brings its own rate.
\item Fixed numbers have no rate and add no term.
\item Put the moment's numbers in last. Find any missing quantity from the link equation.
\item Say the answer with its sign and units: positive grows, negative shrinks.\end{itemize}
\section*{Every level, with answers}
\subsection*{Differentiate the link}
\noindent\textbf{Level 1.} The link is $A=\pi r^2$. All the letters change with time $t$. Which line is the result of differentiating both sides with respect to $t$?\\ \A{(A) $\dfrac{dA}{dt}=2\pi r\,\dfrac{dr}{dt}$} \quad (B) $\dfrac{dA}{dt}=2\pi\,\dfrac{dr}{dt}$ \quad (C) $\dfrac{dA}{dt}=\pi r^2\,\dfrac{dr}{dt}$ \quad (D) $\dfrac{dA}{dt}=2\pi r$\\ \A{Answer: A} \quad Differentiate each side in time. Each changing letter brings its own rate with it, so the answer is $\dfrac{dA}{dt}=2\pi r\,\dfrac{dr}{dt}$.\par\medskip
\noindent\textbf{Level 2.} The link is $x^2+y^2=z^2$. All the letters change with time $t$. Which line is the result of differentiating both sides with respect to $t$?\\ (A) $2x\,\dfrac{dx}{dt}+2y\,\dfrac{dy}{dt}=2z$ \quad (B) $x\,\dfrac{dx}{dt}+y\,\dfrac{dy}{dt}=z\,\dfrac{dz}{dt}$ \quad (C) $2x+2y=2z$ \quad \A{(D) $2x\,\dfrac{dx}{dt}+2y\,\dfrac{dy}{dt}=2z\,\dfrac{dz}{dt}$}\\ \A{Answer: D} \quad Differentiate each side in time. Each changing letter brings its own rate with it, so the answer is $2x\,\dfrac{dx}{dt}+2y\,\dfrac{dy}{dt}=2z\,\dfrac{dz}{dt}$.\par\medskip
\noindent\textbf{Level 3.} The link is $V=\pi r^2h$ (both $r$ and $h$ change). All the letters change with time $t$. Which line is the result of differentiating both sides with respect to $t$?\\ (A) $\dfrac{dV}{dt}=2\pi r\,\dfrac{dr}{dt}+\pi\,\dfrac{dh}{dt}$ \quad (B) $\dfrac{dV}{dt}=2\pi rh\,\dfrac{dr}{dt}$ \quad (C) $\dfrac{dV}{dt}=\pi r^2\,\dfrac{dh}{dt}$ \quad \A{(D) $\dfrac{dV}{dt}=2\pi rh\,\dfrac{dr}{dt}+\pi r^2\,\dfrac{dh}{dt}$}\\ \A{Answer: D} \quad Differentiate each side in time. Each changing letter brings its own rate with it, so the answer is $\dfrac{dV}{dt}=2\pi rh\,\dfrac{dr}{dt}+\pi r^2\,\dfrac{dh}{dt}$.\par\medskip
\subsection*{Solve for the wanted rate}
\noindent\textbf{Level 1.} A circular puddle's radius grows 3 cm every second. Using $\dfrac{dA}{dt}=2\pi r\dfrac{dr}{dt}$, find how fast the area grows when $r=7$ cm. Give it to 2 decimal places.\\ \A{Answer: 131.95 cm\textsuperscript{2}/s} \quad $\dfrac{dA}{dt}=2\pi(7)(3)=131.95$ cm\textsuperscript{2}/s.\par\medskip
\noindent\textbf{Level 2.} A balloon's radius grows 4 cm every second. Using $\dfrac{dV}{dt}=4\pi r^2\dfrac{dr}{dt}$, find how fast the volume grows when $r=8$ cm. Give it to 2 decimal places.\\ \A{Answer: 3216.99 cm\textsuperscript{3}/s} \quad $\dfrac{dV}{dt}=4\pi(8)^2(4)=3216.99$ cm\textsuperscript{3}/s.\par\medskip
\noindent\textbf{Level 3.} A 5 ft ladder slides away from a wall. The foot moves out at 3 ft/s. Using $2x\dfrac{dx}{dt}+2y\dfrac{dy}{dt}=0$, find $\dfrac{dy}{dt}$ when the foot is 3 ft from the wall and the top is 4 ft up. Give it to 2 decimal places.\\ \A{Answer: -2.25 ft/s} \quad $18+8\dfrac{dy}{dt}=0$, so $\dfrac{dy}{dt}=-\dfrac{18}{8}=-2.25$ ft/s. The top is falling, so it is negative.\par\medskip
\subsection*{Say what the rate means}
\noindent\textbf{Level 1.} You solve a related-rates problem and get $\dfrac{dA}{dt}=4.8$ cm\textsuperscript{2}/s. What does that say?\\ \A{(A) The area is growing 4.8 square centimeters per second at that moment.} \quad (B) The area is 4.8 cm\textsuperscript{2} at that moment. \quad (C) The area is growing 4.8 cm per second at that moment. \quad (D) The area is shrinking 4.8 square centimeters per second at that moment.\\ \A{Answer: A} \quad The sign is positive, so growing, and cm\textsuperscript{2}/s is square centimeters per second. The area is growing 4.8 square centimeters per second at that moment.\par\medskip
\noindent\textbf{Level 2.} You solve a related-rates problem and get $\dfrac{dV}{dt}=-2.0$ cm\textsuperscript{3}/s. What does that say?\\ (A) The volume is growing 2.0 cubic centimeters per second at that moment. \quad (B) The volume is shrinking 2.0 cm per second at that moment. \quad \A{(C) The volume is shrinking 2.0 cubic centimeters per second at that moment.} \quad (D) The volume is 2.0 cm\textsuperscript{3} at that moment.\\ \A{Answer: C} \quad The sign is negative, so shrinking, and cm\textsuperscript{3}/s is cubic centimeters per second. The volume is shrinking 2.0 cubic centimeters per second at that moment.\par\medskip
\noindent\textbf{Level 3.} You solve a related-rates problem and get $\dfrac{dA}{dt}=9.0$ ft\textsuperscript{2}/s. What does that say?\\ (A) The area is 9.0 ft\textsuperscript{2} at that moment. \quad (B) The area is shrinking 9.0 square feet per second at that moment. \quad \A{(C) The area is growing 9.0 square feet per second at that moment.} \quad (D) The area is growing 9.0 ft per second at that moment.\\ \A{Answer: C} \quad The sign is positive, so growing, and ft\textsuperscript{2}/s is square feet per second. The area is growing 9.0 square feet per second at that moment.\par\medskip
\subsection*{Fixed numbers have no rate}
\noindent\textbf{Level 1.} A can has a FIXED radius of 4 cm. Its water depth $h$ changes. The link is $V=\pi r^2h$. Which line differentiates it in time?\\ \A{(A) $\dfrac{dV}{dt}=\pi r^2\,\dfrac{dh}{dt}$} \quad (B) $\dfrac{dV}{dt}=0$ \quad (C) $\dfrac{dV}{dt}=\pi r^2$ \quad (D) $\dfrac{dV}{dt}=2\pi rh\,\dfrac{dr}{dt}$\\ \A{Answer: A} \quad Fixed numbers have rate 0, so they add no term. Only the changing letters bring rates: $\dfrac{dV}{dt}=\pi r^2\,\dfrac{dh}{dt}$.\par\medskip
\noindent\textbf{Level 2.} A box has a FIXED length 5 and width 3. Its height $h$ changes. The link is $V=lwh$. Which line differentiates it in time?\\ (A) $\dfrac{dV}{dt}=0$ \quad (B) $\dfrac{dV}{dt}=w h\,\dfrac{dl}{dt}+l h\,\dfrac{dw}{dt}$ \quad (C) $\dfrac{dV}{dt}=lw$ \quad \A{(D) $\dfrac{dV}{dt}=lw\,\dfrac{dh}{dt}$}\\ \A{Answer: D} \quad Fixed numbers have rate 0, so they add no term. Only the changing letters bring rates: $\dfrac{dV}{dt}=lw\,\dfrac{dh}{dt}$.\par\medskip
\noindent\textbf{Level 3.} A plane flies level at a FIXED height of 3 miles. Its distance along the ground is $x$ and its distance from a station is $s$. The link is $x^2+3^2=s^2$. Which line differentiates it in time?\\ (A) $2x=2s$ \quad (B) $2x\,\dfrac{dx}{dt}+6=2s\,\dfrac{ds}{dt}$ \quad \A{(C) $2x\,\dfrac{dx}{dt}=2s\,\dfrac{ds}{dt}$} \quad (D) $2x\,\dfrac{dx}{dt}+2\cdot3\,\dfrac{dy}{dt}=2s\,\dfrac{ds}{dt}$\\ \A{Answer: C} \quad Fixed numbers have rate 0, so they add no term. Only the changing letters bring rates: $2x\,\dfrac{dx}{dt}=2s\,\dfrac{ds}{dt}$.\par\medskip
\section*{Common misconceptions}\begin{tabularx}{\textwidth}{@{}X X@{}}\toprule Students say & Say instead\\\midrule
dA/dt = 2$\pi$r (stop there). & The chain factor dr/dt rides along: dA/dt = 2$\pi$r$\cdot$dr/dt.\\
Put r = 10 in first, then differentiate. & A fixed 10 has rate 0, so everything vanishes. Differentiate first, plug in after.\\
A fixed number like 13 or the radius of a fixed can needs a rate. & Fixed numbers have rate 0 and add no term.\\
A shrinking quantity has a positive answer. & Shrinking means the rate is negative.\\
Product rule: multiply the two rates. & Each term pairs one quantity with the rate of the other.\\
\bottomrule\end{tabularx}
\end{document}